\section{Training, Baselines, and Generation Integrity}
\label{app:training-baselines}

This appendix separates the controlled Wan comparison from the external
CogVideoX benchmark block.  Within Wan, the Matched Control and \methodshort{}
differ only in the identity placed in the student prompt's registered source
slot.  The CogVideoX methods instead run on their own frozen backbone and are
interpreted relative to CogVideoX Original; raw Wan--CogVideoX differences are
descriptive rather than same-backbone method ablations.

\subsection{Single-Factor Matched Design}
\label{app:matched-design}

The main comparison isolates source-slot replacement.  For each mechanism,
the Matched Control conditions the student on the original factual prompt
\(p_i^{\mathrm f}\), whereas \methodshort{} conditions it on the reconstructed
prompt \(p_i^{\mathrm{aug}}\).  Table~\ref{tab:matched-contract} lists the
quantities shared by construction.

\begin{table}[H]
  \centering
  \caption{Registered single-factor contract between the two Wan arms.  A
  shared entry is hash- or value-bound before training.  The arm-specific
  student condition is the only intended treatment difference.}
  \label{tab:matched-contract}
  \small
  \setlength{\tabcolsep}{4pt}
  \begin{tabularx}{\linewidth}{@{}>{\raggedright\arraybackslash}p{0.27\linewidth}>{\raggedright\arraybackslash}X>{\raggedright\arraybackslash}X@{}}
    \toprule
    Quantity & Matched Control & \methodshort{} \\
    \midrule
    Student factual condition
      & Original source in \(p_i^{\mathrm f}\)
      & Assigned bank source in \(p_i^{\mathrm{aug}}\) \\
    Counterfactual media and latent
      & \multicolumn{2}{c}{Shared \(x_i^{\mathrm{cf}}\) and \(z_i^{\mathrm{cf}}\)} \\
    Frozen-teacher condition
      & \multicolumn{2}{c}{Shared \(p_i^{\mathrm{cf}}\) on the same noised latent} \\
    Preservation branch
      & \multicolumn{2}{c}{Same 36 videos, prompts, latents, and embeddings} \\
    Erasure indices and update order
      & \multicolumn{2}{c}{Same active 100 indices and alternating schedule} \\
    Initialization and stochastic draws
      & \multicolumn{2}{c}{Same LoRA initialization, noise, \(\sigma\), and timestep trace} \\
    Optimizer, checkpoint, and inference
      & \multicolumn{2}{c}{Identical registered configuration} \\
    \bottomrule
  \end{tabularx}
\end{table}

Seven equality receipts cover the 178 Matched student embeddings in each
mechanism and require every fresh Matched encoding to equal its base-cache
factual embedding tensor exactly.  Each training run then repeats an A--A
numerical preflight: forward tensors and losses must match exactly, while the
registered gradient tolerances require maximum absolute difference at most
\(10^{-4}\), per-parameter relative \(\ell_2\) difference at most \(0.15\),
and global relative \(\ell_2\) difference at most \(0.05\).  This preflight
checks exact A--A forward/loss repeatability and bounded numerical variation
in gradients; it does not require the
Matched and \methodshort{} gradients or optimized trajectories to coincide
after their student prompts diverge.

\subsection{Optimization and Inference}
\label{app:optimization-details}

Erasure and preservation are separate alternating updates.  On an erasure
step, the student receives the arm-specific factual condition and a noised
source-free target latent, while the adapter-disabled transformer supplies the
online target-prompt prediction on the same latent and timestep.  On a
preservation step, the student and frozen transformer receive the same generic
prompt and noised preservation latent.  The implementation never sums the two
objectives on one minibatch.

\begin{table}[H]
  \centering
  \caption{Common Wan training and formal-inference configuration.  Every one
  of the 18 registered Wan runs uses these values without mechanism-specific
  tuning.}
  \label{tab:wan-optimization-config}
  \small
  \begin{tabularx}{\linewidth}{@{}>{\raggedright\arraybackslash}p{0.36\linewidth}>{\raggedright\arraybackslash}X@{}}
    \toprule
    Field & Registered value \\
    \midrule
    Base model & Wan 2.1 T2V 1.3B Diffusers \\
    Adapter unit & One mechanism-specific LoRA per training run \\
    Training data & 178 erasure bindings; 36 preservation bindings \\
    Active schedule & 200 updates, strictly alternating 100 erasure and 100 preservation \\
    Training seed / precision & 26000 / bfloat16 \\
    LoRA & Rank 16, alpha 16, on \texttt{to\_q}, \texttt{to\_k}, \texttt{to\_v}, and \texttt{to\_out.0} \\
    Optimizer & AdamW, learning rate \(5\!\times\!10^{-5}\), betas \((0.9,0.999)\), weight decay 0.01 \\
    Batch / accumulation / clipping & 1 / 1 / maximum gradient norm 1.0 \\
    Loss weights & \(\lambda_{\mathrm T}=4\), \(\lambda_{\mathrm P}=4\) \\
    Checkpoint policy & Step 200 only \\
    Formal inference & LoRA scale 1.25; guidance 5; 25 denoising steps \\
    Video specification & 49 frames, 8 fps, \(832\!\times\!480\), bfloat16 \\
    \bottomrule
  \end{tabularx}
\end{table}

All 18 runs produced an eligible step-200 checkpoint.  Their receipts record
100 erasure and 100 preservation updates, a passed finite-parameter check, and
the precommitted initialization, schedule, and noise-generator digests.  The
formal generation runner verified the corresponding receipt and checkpoint
hashes before producing any output.  These are 18 mechanism/control
operators, not one jointly trained multi-mechanism adapter.

\subsection{Identification Training Arms}
\label{app:identification-arms}

The identification study varies the surface realization and causal or
noncausal placement of source vocabulary without changing the target, teacher,
preservation branch, objective, schedule, or checkpoint policy.  The main
design trains seven Matched and seven
\methodshort{} adapters.  Water Impact and Brittle Fracture additionally
receive a Generic Paraphrase and a Bystander Token adapter, yielding 18 Wan
runs in total.

\begin{table}[H]
  \centering
  \caption{Student-prompt intervention in each Wan training arm.  The two
  identification controls are restricted to Water Impact and Brittle
  Fracture.}
  \label{tab:identification-arms}
  \small
  \setlength{\tabcolsep}{4pt}
  \begin{tabularx}{\linewidth}{@{}>{\raggedright\arraybackslash}p{0.23\linewidth}>{\raggedright\arraybackslash}X>{\raggedright\arraybackslash}p{0.22\linewidth}@{}}
    \toprule
    Arm & Student factual condition & Runs \\
    \midrule
    Matched Control
      & Original causal-source identity and registered wording
      & Seven, one per mechanism \\
    Generic Paraphrase
      & Original causal-source identity with the opposite registered direct/natural wording
      & Two: Water, Fracture \\
    Bystander Token
      & Original causal source retained; the \methodshort{}-assigned identity appears motionless in a far-background noncausal clause
      & Two: Water, Fracture \\
    \methodshort{}
      & \methodshort{}-assigned identity occupies the causal-source slot
      & Seven, one per mechanism \\
    \bottomrule
  \end{tabularx}
\end{table}

The Bystander Token construction matches the assigned-identity frequency
distribution of \methodshort{} exactly while leaving the true causal source in
place.  Generic Paraphrase instead preserves the original source and changes
only its registered surface realization.  Thus these arms test paraphrase
exposure and noun-frequency placement; they are not component-removal
ablations of the losses in Sec.~\ref{subsec:training-objectives}.  The frozen
identification subset contains 24 cases from each of Water and Fracture.  The
Matched and \methodshort{} outputs are reused from the main evaluation, while
Generic Paraphrase and Bystander Token add \(2\times2\times24=96\) videos.

\subsection{External Baseline Implementations}
\label{app:baseline-implementations}

The external block uses one frozen CogVideoX-2B model inventory and supplies
the same positive prompt and case seed to every compatible interface.
CogVideoX Original receives no erasure concept.  Negative Prompt, VideoEraser,
and SAFREE receive the registered mechanism name through their respective
concept-string interfaces.  T2VUnlearning-adapted instead loads a separately
trained operator for each mechanism and receives no concept string at
inference.  Table~\ref{tab:external-baselines} distinguishes official
pipelines from the internal adaptation required when matching public training
code and mechanism checkpoints were unavailable.

\begin{table}[H]
  \centering
  \caption{CogVideoX benchmark implementations and fidelity boundaries.}
  \label{tab:external-baselines}
  \footnotesize
  \setlength{\tabcolsep}{4pt}
  \begin{tabularx}{\linewidth}{@{}>{\raggedright\arraybackslash}p{0.25\linewidth}>{\raggedright\arraybackslash}X@{}}
    \toprule
    Display name & Registered implementation \\
    \midrule
    CogVideoX Original
      & Unmodified positive-prompt generation on CogVideoX-2B. \\
    \negativeprompt
      & The mechanism name is supplied through the native negative-prompt interface. \\
    \videoeraser
      & Unmodified official CogVideoX SPEA/ARNG pipeline snapshot; the mechanism name is passed through its unsafe-concept interface. \\
    \tvtwounlearning
      & Paper-guided internal CogVideoX adaptation using negatively guided velocity matching and an unmasked adapter-residual magnitude penalty.  It does not claim parity with the unreleased QK/target-token mask or prompt-augmentation training procedure. \\
    \safree
      & Official CogVideoX pipeline snapshot with the mechanism name inserted as one registered in-memory concept entry; this interface runs in fp32. \\
    \bottomrule
  \end{tabularx}
\end{table}

The VideoEraser snapshot is frozen at commit
\texttt{ba19cceb561d}, with pipeline SHA-256 prefix
\texttt{bd9e4052740e}.  The SAFREE snapshot is frozen at commit
\texttt{b8b2c3fa9d7f}, with pipeline SHA-256 prefix
\texttt{55185f5972ec}.  The complete revisions and digests are retained in the
frozen baseline registry.
For \tvtwounlearning, 36 prompts per mechanism are selected by a frozen
hash rule and train rank-128 residual adapters on every CogVideoX
\texttt{attn1} block.  Its precommitted configuration uses seed 12000,
learning rate \(10^{-4}\), negative-velocity scale 7, residual-magnitude
weight 1,
no preservation term, and 100 updates; step 100 is the only eligible
checkpoint.  Seven such checkpoints were frozen before formal generation.

\subsection{Runtime and Generation Accounting}
\label{app:generation-accounting}

The frozen execution environment uses CPython 3.11.15, PyTorch 2.6.0 with
CUDA 12.4, Diffusers 0.33.1, Transformers 4.51.3, and four NVIDIA
A100-SXM4-80GB devices.  Table~\ref{tab:generation-config} records the two
backbone-specific generation configurations.  Matching holds within each
block; the different backbone, resolution, denoising budget, guidance, and
SAFREE precision prevent raw cross-block scores from being interpreted as a
pure method comparison.

\begin{table}[H]
  \centering
  \caption{Backbone-specific formal generation settings.}
  \label{tab:generation-config}
  \small
  \setlength{\tabcolsep}{4pt}
  \begin{tabularx}{\linewidth}{@{}>{\raggedright\arraybackslash}p{0.27\linewidth}>{\raggedright\arraybackslash}X>{\raggedright\arraybackslash}X@{}}
    \toprule
    Field & Wan block & CogVideoX block \\
    \midrule
    Backbone & Wan 2.1 T2V 1.3B & CogVideoX-2B \\
    Resolution & \(832\!\times\!480\) & \(720\!\times\!480\) \\
    Frames / fps & 49 / 8 & 49 / 8 \\
    Denoising steps & 25 & 50 \\
    Guidance scale & 5 & 6 \\
    Scheduler & Frozen Wan pipeline default & CogVideoX DPM, trailing spacing, dynamic CFG \\
    Precision & bfloat16 & Mixed: core backbone in bfloat16; T2V residual adapters and SAFREE in fp32 \\
    Pairing & Same case prompt and seed within Wan & Same case prompt and seed within CogVideoX \\
    \bottomrule
  \end{tabularx}
\end{table}

Every generated file was decoded before publication into the anonymous review
package and checked for one video stream, no audio, 49 frames, 8 fps, the
registered resolution, and a byte hash matching its generation receipt.
Table~\ref{tab:generation-accounting} reconciles the complete package.

\begin{table}[H]
  \centering
  \caption{Frozen formal-generation accounting.  The trained-Wan block
  contains 588 main Matched/\methodshort{} videos and 96 additional
  identification-control videos.}
  \label{tab:generation-accounting}
  \small
  \begin{tabular}{@{}lrrr@{}}
    \toprule
    Block & Jobs & Expected & Validated \\
    \midrule
    Wan Original & 7 & 294 & 294 \\
    Trained Wan & 18 & 684 & 684 \\
    CogVideoX core four & 28 & 1,176 & 1,176 \\
    SAFREE-CogVideoX & 7 & 294 & 294 \\
    \midrule
    Total & 60 & 2,448 & 2,448 \\
    \bottomrule
  \end{tabular}
\end{table}

The 2,448 outputs consist of 2,352 main-study videos---294 cases under each
of eight streams---and the 96 identification additions.  The four generation
aggregates report complete expected and validated counts, and the package
receipt confirms exact media and blinding validation.  Its SHA-256 prefix is
\texttt{63933393ea04}.  The final baseline registry is frozen, authorizes
formal generation, and has SHA-256 prefix \texttt{8a9d92b86bd1}; the full
digests are contained in the frozen package receipt and baseline registry.
Generation completion establishes inventory integrity, not a same-backbone
comparison between Wan and CogVideoX or an official-training claim for
\tvtwounlearning.
